\subsubsection*{引理 5.3.1.}
设 $K$ 是剩余域为 $\mathbb F_q$ 的局部域, $G$ 是其 Galois 群, $I$ 是惯性群,
$\rho$ 是 $G$ 在 $n$ 维 $E$ 向量空间 $V$ 上的连续表示 (参见 5.0).
用 $V^I$ 和 $V_I$ 分别表示 $I$ 的不变元空间和余不变元空间, 视为 $G/I$ 模.
置 $r=\dim_E V^I$, 并设 $\alpha_1,\ldots,\alpha_r$ 为 $F$ 在 $V^I$ 上的特征值.
则存在整数 $m_i\geq0$, 使 $F$ 在 $V_I$ 上的特征值为
\[
\alpha_iq^{m_i}\qquad(1\leq i\leq r).
\]

\medskip
\noindent\emph{证明.}
设 $I'$ 为 $I$ 中阶与 $\ell$ 互素的最大子群. 群 $\rho(I')$ 有限,
因为它是 $\operatorname{GL}(V)$ 的阶与 $\ell$ 互素的紧子群. 故
\[
V_I\simeq\bigl(V^{\rho(I')}\bigr)_I.
\]
因此将 $V$ 换成 $V^{I'}$ 后, 可设 $I'$ 在 $V$ 上平凡作用. 作为 $G/I$ 模有
\[
I/I'\simeq\mathbb Z_\ell(1).
\]
设 $T$ 是 $I/I'$ 的拓扑生成元. 对每个 $\alpha$, 设 $V^{(\alpha)}$
为 $T$ 在 $V$ 上对应特征值 $\alpha$ 的广义特征空间. 则
\[
V=\bigoplus_\alpha V^{(\alpha)},
\]
且分解
\[
V=V^{(1)}\oplus\bigoplus_{\alpha\ne1}V^{(\alpha)}
\]
在 $G$ 下稳定并与对偶相容. 因此可将 $V$ 换成 $V^{(1)}$, 并设 $I$ 幺幂作用.

\subsubsection*{引理 5.3.2.}
设 $V$ 是特征 $0$ 的域 $k$ 上的有限维向量空间, $I$ 是单参数加法群,
$\rho:I\longrightarrow\operatorname{GL}(V)$ 是表示, $N$ 是 $I$ 的无穷小生成元. 置
\[
V^{I,k}=V^I\cap\operatorname{Im}(N^k)
\]
及
\[
V_{I,k}=\operatorname{Im}\!\left(\operatorname{Ker}(N^{k+1})
\longrightarrow V_I\right).
\]
则 $N^k$ 诱导同构
\[
N^k:V_{I,k}/V_{I,k-1}\xrightarrow{\sim}
V^{I,k}/V^{I,k+1},
\]
并定义典范同构
\[
V_{I,k}/V_{I,k-1}\otimes I^{\otimes k}
\simeq V^{I,k}/V^{I,k+1}.
\]

\medskip
\noindent\emph{证明.}
通过 $N$, $I$ 的表示 $V$ 可等同于被 $T$ 的某次幂化为零的 $k[T]$ 模.
只需在 $V$ 对应不可分解 $k[T]$ 模 $k[T]/T^k$ 时验证引理~5.3.2;
此时断言立即成立.

现完成引理~5.3.1 的证明. 引理~5.3.2 分别给出 $V_I$ 和 $V^I$ 的递增滤过 $V_{I,k}$
与递减滤过 $V^{I,k}$, 以及 $G/I$ 模同构
\[
V^{I,k}/V^{I,k+1}\simeq
\bigl(V_{I,k}/V_{I,k-1}\bigr)(k).
\]
断言得证.

\medskip
\noindent\emph{底本注.}
照录法文层用 $V$ 的维数 $n$ 枚举 $V^I$ 上的特征值, 将引理引用为“3.2”,
并称 $k[T]/T^k$ 不可约. 校订法文层和英文层采用 $r=\dim_E V^I$、准确编号 5.3.2,
以及数学上正确的术语\emph{不可分解}.
